\documentclass[12pt]{article}
\usepackage[margin=1in]{geometry}
\usepackage[T1]{fontenc}
\usepackage[utf8]{inputenc}
\usepackage{mathptmx}
\setlength{\parindent}{0pt}
\setlength{\parskip}{1em}

\begin{document}

\thispagestyle{empty}

\begin{center}
    {\LARGE\bfseries Executive Summary}\\[1.2em]
    {\large Ephemeral Interfaces: Design Patterns for Task-Scoped Generative UI Components}\\[0.8em]
    {Daniel Kumlin}
\end{center}

Generative UI has become a broader topic in software design because it promises interfaces that adapt to user goals instead of remaining static. Yet this promise collides with modern software use, where users still expect stable experiences, predictable controls, and durable mental models. This thesis investigates whether ephemeral, task-scoped generative UI components can offer a bounded alternative to full interface regeneration by appearing only when locally useful and disappearing without destabilising the host application.

The thesis develops and tests a framework for ephemeral support through a web-based research artifact. The main individual contributions are the framework operationalisation, a constrained component-specification model for temporary support, a validation pipeline for generated interface specifications, a bounded ephemeral overlay approach that preserves interface continuity, and an interactive survey artifact for empirical evaluation. The artifact was tested in a within-subjects user study comparing a baseline interface with an ephemeral condition on a bounded slide-reordering task.

The results show a mixed outcome. The ephemeral condition showed a descriptively higher task completion rate and significantly higher helpfulness ratings, while perceived control remained stable, but it did not improve completion time or overall preference and significantly increased interaction count and perceived intrusiveness. These findings suggest that ephemeral generative support is most credible not as a general replacement for existing interfaces, but as a bounded design strategy for bounded, low-consequence subtasks within broader workflows, especially where lightweight contextual guidance may be useful and interruption costs remain manageable. The thesis therefore contributes a framework, design guidelines, and an empirical basis for assessing when temporary generative assistance can add value without disrupting the user experience.
\end{document}
